% ======================================================================
%  PARTIE 6 -- Expérimentations et résultats                 (~4 min)
% ======================================================================
\section{Expérimentations et résultats}

% ----------------------------------------------------------------------
% À DIRE (~30 s)
%  - Deux robots terrestres AgileX :
%      * Scout Mini : compact, 4 roues motrices indépendantes, 10 kg de
%        charge utile, 1,5 m/s max, ~3 h d'autonomie, pentes jusqu'à 30° ;
%      * Ranger : 4 roues indépendantes, omnidirectionnel, jusqu'à 150 kg,
%        2,6 m/s max, ~4 h d'autonomie, batteries remplaçables à chaud.
%  - Tous deux équipés du kit R&D de Génération Robots : LiDAR, caméra de
%    profondeur Intel RealSense D435, GPS RTK, IMU, calculateur embarqué
%    avec les pilotes ROS 2.
%  - Acquisitions enregistrées (bags) puis rejouées, dans la même zone
%    (près du bâtiment ESPRIT) -> carte a priori disponible pour les deux :
%      * Ranger : bag d'une acquisition multi-robots existante (drone +
%        Ranger), 8 min, que je n'ai pas réalisée moi-même ;
%      * Scout : acquisition que j'ai pilotée (avec Gérald Dherbomez,
%        plateforme PRETIL) : bag court, facile à rejouer, et surtout une
%        vraie boucle fermée pour tester la loop closure.
% ----------------------------------------------------------------------
\begin{frame}{Plateformes et acquisitions}
  \begin{columns}[c]
    \begin{column}{0.3\textwidth}
      \centering
      \adjincludegraphics[width=\linewidth, trim={{0.08\width} {0.1\height} {0.09\width} {0.22\height}}, clip]{scout.png}\\[2pt]
      \textbf{Scout Mini}\\
      {\footnotesize\color{textgray} 1,5 m/s \textbullet\ 10 kg \textbullet\ 3 h}
    \end{column}
    \begin{column}{0.3\textwidth}
      \centering
      \adjincludegraphics[width=\linewidth, trim={{0.12\width} {0.08\height} {0.15\width} {0.18\height}}, clip]{ranger.png}\\[2pt]
      \textbf{Ranger}\\
      {\footnotesize\color{textgray} 2,6 m/s \textbullet\ 150 kg \textbullet\ 4 h}
    \end{column}
    \begin{column}{0.34\textwidth}
      \begin{tcolorbox}[carte, title=Kit R\&D (Génération Robots)]
        LiDAR 3D \textbullet\ IMU\\
        Caméra RealSense D435\\
        GPS RTK\\
        PC embarqué + \alert{ROS 2}
      \end{tcolorbox}
    \end{column}
  \end{columns}
  \vspace{0.35cm}
  {\small
  \begin{tabular}{@{}l@{\hspace{0.4cm}}l@{}}
    \textbf{Ranger} & bag existant (drone + Ranger), \alert{8 min} \\
    \textbf{Scout Mini} & acquisition pilotée, \alert{boucle fermée} sur le parking \\
  \end{tabular}\\[3pt]
  $\Rightarrow$ même zone (bâtiment \alert{ESPRIT}) : carte a priori disponible pour les deux\par}
\end{frame}

% ----------------------------------------------------------------------
% À DIRE (~40 s)
%  - Résultats sur le bag du Ranger (8 min), affichés sur la carte a
%    priori (qui sert ici uniquement de fond de carte).
%  - L'information sémantique des cases est cohérente avec la réalité : on
%    retrouve les bâtiments et les arbres, visibles aussi sur la carte de
%    LIO (trajectoire en vert).
%  - La carte 3D permet de mieux visualiser bâtiments et arbres.
%  - Limite : un écart de position avec la carte a priori, dû au manque de
%    précision de la position initiale donnée par le GPS (~5 m d'erreur) ;
%    le magnétomètre s'est révélé inutilisable, l'orientation initiale a
%    donc été recalée visuellement.
% ----------------------------------------------------------------------
\begin{frame}{Résultats : Ranger}
  \centering
  \begin{tabular}{@{}c@{\hspace{0.3cm}}c@{\hspace{0.3cm}}c@{}}
    \adjincludegraphics[height=3.6cm, trim={{0.22\width} {0.38\height} {0.2\width} {0.1\height}}, clip]{ranger_carte_zoom.png} &
    \adjincludegraphics[height=3.6cm, trim={{0.235\width} {0.13\height} {0.18\width} {0.35\height}}, clip]{ranger_liosam.png} &
    \adjincludegraphics[height=3.6cm, trim={{0.0\width} {0.08\height} {0.0\width} {0.0\height}}, clip]{ranger_carte3d.png} \\
    \legende{\shortstack{Carte 2D (pas de 1 m)\\ sur la carte a priori}} &
    \legende{Carte et trajectoire LIO} &
    \legende{Carte sémantique 3D} \\
  \end{tabular}\par
  \vspace{0.2cm}
  {\small \ok\ sémantique \alert{cohérente} (bâtiments, arbres)
   \qquad \ko\ \alert{décalage} : GPS initial ($\sim$5 m)\par}
\end{frame}

% ----------------------------------------------------------------------
% À DIRE (~50 s)
%  - Bag du Scout : un tour complet du parking devant ESPRIT, à vitesse de
%    marche -> vraie loop closure en conditions réelles.
%  - Sans loop closure : la carte est dédoublée, à cause de la dérive de
%    la trajectoire estimée par LIO. À un moment la dérive devient très
%    forte : LIO retrouve ensuite une trajectoire cohérente mais avec une
%    pente très exagérée, alors que le robot est resté à plat.
%    Hypothèse : données IMU brutes, non filtrées (pas de dérive de ce
%    type avec le Ranger).
%  - Avec loop closure : carte nettement plus cohérente avec le chemin
%    réel. L'optimiseur ne corrige pas que la dernière pose : il
%    redistribue la correction sur toute la trajectoire -> c'est
%    exactement ce que fait le recalcul complet des points.
%  - Remarque : avec loop closure, même le DÉBUT de la carte change
%    légèrement d'orientation par rapport au cap initial renseigné à la
%    main, car la pose de départ n'est pas figée dans l'optimisation.
%  - (Slide de secours en annexe : bag lu à partir de 60 s.)
% ----------------------------------------------------------------------
\begin{frame}{Résultats : Scout Mini et loop closure}
  \centering
  % images centrées verticalement (\raisebox) pour aligner les étiquettes de ligne
  \begin{tabular}{@{}c@{\hspace{0.25cm}}c@{\hspace{0.5cm}}c@{}}
    & \textbf{Sans loop closure} & \textbf{Avec loop closure} \\[3pt]
    \rotatebox[origin=c]{90}{\scriptsize\color{textgray} carte sémantique} &
    \raisebox{-0.5\height}{\adjincludegraphics[height=2.45cm, trim={{0.521\width} {0.300\height} {0.122\width} {0.455\height}}, clip]{scout_sansloop.png}} &
    \raisebox{-0.5\height}{\adjincludegraphics[height=2.45cm, trim={{0.503\width} {0.304\height} {0.140\width} {0.440\height}}, clip]{scoot_loop.png}} \\[1.3cm]
    \rotatebox[origin=c]{90}{\scriptsize\color{textgray} trajectoire LIO} &
    \raisebox{-0.5\height}{\includegraphics[height=2.45cm]{scoot_liosam_sansloop.png}} &
    \raisebox{-0.5\height}{\includegraphics[height=2.45cm]{scoot_liosam_loop.png}} \\[1.35cm]
    & {\footnotesize \ko\ \alert{dérive} : carte dédoublée, pente exagérée}
    & {\footnotesize \ok\ carte \alert{cohérente}, correction globale} \\
  \end{tabular}
\end{frame}

% ----------------------------------------------------------------------
% À DIRE (~40 s)
%  - Test multi-robots avec les deux bags (Ranger et Scout) rejoués en
%    même temps : ce n'est donc pas du multi-robots temps réel (les
%    acquisitions n'ont pas été faites simultanément).
%  - Mais cela valide les bases : les deux cartes sont agrégées dans une
%    carte commune, un seul fichier GeoJSON, sans conflit de topics ni de
%    repères TF.
%  - Limite : si une même case est vue par les deux robots, les deux
%    informations sont gardées -> deux lignes pour la même case dans le
%    GeoJSON -> il faudra une fusion (bayésienne).
%  - La loop closure fonctionne aussi en multi-robots, mais avec les deux
%    bags + deux frameworks sur la même machine, les pertes de données
%    sont encore plus importantes.
% ----------------------------------------------------------------------
\begin{frame}{Résultats : multi-robots}
  \begin{columns}[c]
    \begin{column}{0.53\textwidth}
      \centering
      \begin{tikzpicture}
        % carte complète (hauteur 5 cm -> largeur 3,855 cm) et zoom sur les deux zones
        \node[inner sep=0pt, anchor=south west] (full) at (0,0) {\includegraphics[height=5cm]{multi.png}};
        \node[inner sep=0pt, anchor=south west] (zoom) at (4.15,0)
          {\adjincludegraphics[height=5cm, trim={{0.565\width} {0.291\height} {0.106\width} {0.304\height}}, clip]{multi.png}};
        \draw[accent, thick] (2.178,1.455) rectangle (3.446,3.48);
        \draw[accent, thin] (3.446,3.48) -- (zoom.north west);
        \draw[accent, thin] (3.446,1.455) -- (zoom.south west);
        \draw[accent, thick] (zoom.south west) rectangle (zoom.north east);
        \node[anchor=north, font=\scriptsize, text=textgray] at (full.south) {carte commune};
        \node[anchor=north, font=\scriptsize, text=textgray] at (zoom.south) {zoom};
        \node[anchor=east, fill=white, inner sep=1pt, font=\tiny\bfseries, text=primaryblue] at ([xshift=-2pt, yshift=-1.4cm]zoom.north east) {Ranger};
        \node[anchor=east, fill=white, inner sep=1pt, font=\tiny\bfseries, text=primaryblue] at ([xshift=-2pt, yshift=1.1cm]zoom.south east) {Scout};
      \end{tikzpicture}
    \end{column}
    \begin{column}{0.45\textwidth}
      \begin{itemize}
        \item 2 bags rejoués \alert{ensemble}\\ {\footnotesize\color{textgray} (pas de temps réel)}
        \item[\ok] \alert{un seul GeoJSON} commun
        \item[\ok] pas de conflit de topics ni de TF
        \item[\ko] case vue 2 fois $\rightarrow$ 2 entrées\\ $\Rightarrow$ \alert{fusion bayésienne} à faire
        \item[\ko] loop closure + 2 robots sur 1 PC\\ $\rightarrow$ pertes de données
      \end{itemize}
    \end{column}
  \end{columns}
\end{frame}
